Haim--Kislev and Ostrover disproved Viterbo's conjecture with a Lagrangian
product of two regular pentagons, the Haim--Kislev--Ostrover (HKO) body.
We prove that this body locally maximizes the systolic ratio among
polytopes with exactly ten facets. Among bodies sufficiently near HKO,
equality in this systolic-ratio bound holds precisely for its translations,
positive dilations and linear symplectic images. We also prove an explicit capacity formula and the sharp
bound $(3+\sqrt5)/5$ for products of two independently linearly deformed
regular pentagons. Thus neither sufficiently small ten-facet perturbations
nor arbitrary deformations within this product family improve the HKO ratio.

The placement of product planes supplies another source of variation:
every symplectic product of two full-dimensional planar convex bodies admits
a capacity-increasing rigid rotation, while the Lagrangian endpoint can
have smaller capacity. The proofs and computations use finite capacity
formulas, their first-order variation, and an implementation of
Haim--Kislev's quadratic program. A flow-graph algorithm using the local
passage geometry of Chaidez and Hutchings gives a second method for
capacity computation under explicit regularity assumptions.

In broader sampled populations of four-dimensional convex polytopes,
symplectic ridge area has a strong inverse association with the computed
systolic ratio. Selecting bodies by geometric measurements raises mean
ratios in fresh samples, but the ridge association weakens among low-ridge
or high-ratio subsets. Area-preserving factor deformations that balance
vertex covariances raise ratios in most tested products and weaken their
ridge--ratio association, supporting a contribution from affine factor
shape to that association in the tested sample. This does not give a
universal ordering: exact relative-rotation laws for regular pentagon and
equilateral-triangle products order ridge area and systolic ratio in
opposite ways. None of the exploratory searches produced a new numerical
ratio above one, although retaining previously active capacity branches
improved median best ratios over literal gradient ascent in a matched-start
comparison.
